\section{创业与产业判断：价值优先、垂直突破、效率回归}

\subsection{创业切入点：基础模型之外的价值密度}
在 Q\&A 中，讲者强调从零训练新基础模型仍有空间，但门槛极高，需要资本、人才与算力协同。更现实且机会更大的方向，是结合行业知识在企业场景做垂直化产品，把通用模型能力转化为可交付价值。

\begin{importantbox}{创业建议（工程视角）}
把资源集中在 ``高价值垂直问题''，而不是盲目复制通用模型路线。差异化来自行业数据、流程重构与系统集成深度。
\end{importantbox}

\subsection{中心化与边缘化的动态平衡}
基础训练通常受益于高密度中心集群，但产品推理还受隐私、交互延迟、网络成本和离线可用性约束。本小节不把两者视为互斥路线，而是分析模型大小、任务阶段和数据边界如何决定规划、检索与执行分别放在哪里。
讲者提出一个值得长期追踪的张力：一方面，大规模中心化集群持续推动上限；另一方面，机器人、手机端智能、CDN/电信边缘节点推动高效模型和分布式推理需求。

\begin{knowledgebox}{探索（Exploration）与利用（Exploitation）}
\begin{itemize}
    \item 探索：追求更强模型上限，可容忍更高成本；
    \item 利用：追求可部署效率，强调延迟、功耗与单位成本；
    \item 工业系统通常需要两条路线并行，不存在单一最优策略。
\end{itemize}
\end{knowledgebox}

\subsection{价值创造与风险边界}
讲者对产业持审慎乐观态度：只要系统持续创造用户价值，成本曲线可以被工程优化逐步消化。但若缺乏真实价值，系统最终会成为 ``成本驱动的纸牌屋''。

\begin{warningbox}{价值缺失是最大结构性风险}
如果产品无法在搜索、编码、生产力或企业流程中形成稳定增益，任何技术叙事都难以持续。系统设计的终点必须回到业务价值闭环。
\end{warningbox}

\subsection{本章小结}
未来几年会同时出现 ``超大中心集群'' 与 ``高效边缘推理'' 两条路线。创业公司更适合在垂直场景中把通用能力转化为可量化价值。

